%! Introducing Statistics: Do the Data We Collected Tell the Truth? — Unit 3, Lesson 3.1 — Guided Notes & Practice (Answer Key)
% MIRROR of ../notes/main.tex: -key in place of -boxes, \termblank -> \termans, \blank -> \ans, and the
% fourth argument of every \pcell filled. Nothing else differs — keep it that way.
% THE in-class document, in the MAIN IDEAS / NOTES shape (see LESSON_SHAPE.md §1), at the
% COMPACT budget: the vocabulary box, then ONE guidednotes table of FOUR rows (three instruction
% rows + Guided Practice), 14 problems, 2 pages.
%   I DO   : the teacher fills the definition lines and works the FIRST problem of each grid
%            (1, 4, 8).
%   WE DO  : the class works the rest; the last row, Guided Practice, is worked together.
%   YOU DO : nothing on this page — ../homework/ IS the individual practice.
% THE TRAP is problem 6 (the skate-shop poll's 88% reads as "most people want it").
% THE CRUX is problem 10, in the last instruction row: "2,400 responses is 30 times as many as
% 80, so the 85% is more trustworthy" — size does not repair a non-chance method (EK VAR-1.E.1).
%
% THE CONTEXTS (instruction, guided practice, and homework each use a different one)
%  Row 1 — three one-line methods: radio call-in, a coach's varsity players, 60 random IDs.
%  Row 2 — a city skate park: skate-shop web poll (900 votes, 88%), library on a Tuesday
%          morning (150 people, 34%), 400 random households from the city list (52%).
%  Row 3 — the warm-up's later-start poll comes back: news-site poll, 2,400 responses, 85% yes;
%          80 students drawn at random from the full roster, 58% yes.
%  Row 4 (GP) — a fitness app: step-challenge leaderboard (12,000 users, mean 11,200 steps/day)
%          vs 500 accounts drawn at random from all users (mean 6,900 steps/day).
%
% PAGE LOCKSTEP with notes_key/: \termblank -> \termans, \blank -> \ans, 4th arg of \pcell.
% TABLE MECHANICS: inside a cell `\\` ENDS THE ROW — use \par. Long rows are split into
% sub-rows (a bare `\\` then `& ...`) so the table can break between them.
% NO SKETCH QUESTIONS: every display is a given table, read and compared.
\documentclass[12pt]{article}
\usepackage{apstats-article}
\usepackage{apstats-key}   % pulls in -boxes; do NOT also load -boxes

\begin{document}

\pageheader{Unit 3, Lesson 3.1}{Guided Notes \& Practice --- Answer Key}
% NAMESTRIP: no \namedateperiod here — the name row lives on the cover only.

\vspace{2pt}

\boxguard[16]
\begin{vocabbox}
\small
Fill in each term \textbf{as we name it} in the notes below.
\par
\termans{Data collection method}{how the individuals in the data were chosen and measured.}
\par
\termans{Chance-based method}{a random process decides who is in the data.}
\par
\termans{Non-chance method}{people decide who is in the data --- the researcher or the respondents.}
\par
\termans{Bias}{a method that consistently overestimates or underestimates the truth.}
\par
\end{vocabbox}

\small
\begin{guidednotes}
% ── ROW 1: who decided who is in the data? ───────────────────────────────────
\mainidea[Before you trust it]{Ask how} &
The \textbf{data collection method} is \ans{how} the individuals in the data were chosen and
measured. In a \textbf{chance-based method}, a \ans{random} process decides who is in the data.
In a \textbf{non-chance method}, \ans{people} decide --- the researcher picks whoever is
convenient, or the respondents volunteer themselves.

\vspace{3pt}
\stepnum{1}\ Who is in the data?\par
\stepnum{2}\ Who decided they were in it --- \ans{chance} or people?\par
\stepnum{3}\ Who is left out, and who is over-represented?\par
\stepnum{4}\ So can the result be trusted for \emph{everyone} we want to know about?
\notesprompt{Who decided who is in the data: chance or people?}
\begin{probgrid}{|Y|Y|Y|}
\pcell{1}{A radio host asks listeners to call in and vote on a new stadium.}{1.3cm}{People: the callers chose themselves.} &
\pcell{2}{A coach asks the varsity players if the school should spend more on sports.}{1.3cm}{People: the coach chose.} &
\pcell{3}{A computer randomly picks 60 ID numbers from the full student roster.}{1.3cm}{Chance: a random pick.} \\ \hline
\end{probgrid}
\\ \hline
% ── ROW 2: the method decides whose voice is in the data — THE TRAP ──────────
\mainidea[The method decides]{Who is missing?} &
When people decide who is in the data, some groups are over-represented and others are left
out, so the result is pushed the same way every time. A method has \textbf{bias} if it
\ans{consistently} overestimates or underestimates the truth.
\\
& The city council wants to know what percent of residents support building a skate park.
\par\vspace{3pt}
\begin{minipage}{\linewidth}\centering\footnotesize
\setlength{\tabcolsep}{4pt}%
\renewcommand{\arraystretch}{1.2}
\begin{tabular}{@{} c p{0.58\linewidth} r @{}}
\rowcolor{sky}
& \textbf{Method} & \textbf{Support} \\
\hline
A & Online poll on a local skate shop's website: 900 votes & 88\% \\
B & 150 people asked at the public library, Tuesday at 10 a.m. & 34\% \\
C & 400 addresses picked at random from the city's list of all households & 52\% \\
\end{tabular}
\end{minipage}
\\
& \notesprompt{Use the table.}
\begin{probgrid}{|Y|Y|}
\pcell{4}{Who is likely over-represented in Method A?}{1.5cm}{Skaters and skate-shop fans --- people who want the park.} &
\pcell{5}{Who is likely missing from Method B?}{1.5cm}{People at work or school on a weekday morning.} \\ \hline
\pcell{6}{Is Method A's 88\% likely too high or too low? Why?}{1.5cm}{Too high: fans of the shop chose to vote.} &
\pcell{7}{Which result would you report to the council? Why?}{1.5cm}{C's 52\%: chance picked the households.} \\ \hline
\end{probgrid}
\\ \hline
% ── ROW 3: size does not fix the method — THE CRUX ───────────────────────────
\mainidea[Big is not the same as]{Sample size} &
The warm-up's poll is back. A news site asked, \emph{``Should the school day start later?''}
2{,}400 students responded and \textbf{85\%} said yes. The district then had a computer pick
\textbf{80} students at random from the roster of all students: \textbf{58\%} said yes.
\notesprompt{Compare the two results.}
\begin{probgrid}{|Y|Y|Y|}
\pcell{8}{Which result came from a chance-based method?}{2.0cm}{The 80 picked at random (58\%).} &
\pcell{9}{Why might the news site's 85\% be too high?}{2.0cm}{Students who want a later start chose to answer.} &
\pcell{10}{``2{,}400 is 30 times as many as 80, so the 85\% is more trustworthy.'' Agree? Why?}{2.0cm}{No: volunteers are still biased. Size does not fix the method.} \\ \hline
\end{probgrid}
\\
& \fcolorbox{navy}{sky}{\parbox{\dimexpr\linewidth-2\fboxsep-2\fboxrule\relax}{\centering
\textbf{A bigger sample from a non-chance method is still \ans{biased}.}}}
\\ \hline
% ── ROW 4: guided practice — WE DO, together, fresh context ──────────────────
\mainidea[Guided practice]{The Step Challenge} &
A fitness app wants the average number of steps per day for \emph{all} of its users.
\par\vspace{3pt}
\begin{minipage}{\linewidth}\centering\footnotesize
\setlength{\tabcolsep}{4pt}%
\renewcommand{\arraystretch}{1.2}
\begin{tabular}{@{} p{0.56\linewidth} r @{}}
\rowcolor{sky}
\textbf{Method} & \textbf{Mean steps/day} \\
\hline
Leaderboard: the 12{,}000 users who joined the app's step challenge & 11{,}200 \\
Random: 500 accounts picked at random from all users & 6{,}900 \\
\end{tabular}
\end{minipage}
\notesprompt{We work these together.}
\begin{probgrid}{|Y|Y|}
\pcell{11}{Which method used chance? Who decided in the other one?}{1.5cm}{The random 500. Users chose to join the leaderboard.} &
\pcell{12}{Who is over-represented on the leaderboard?}{1.5cm}{Very active users who like to compete.} \\ \hline
\pcell{13}{Is 11{,}200 likely too high or too low for all users? Why?}{1.5cm}{Too high: the most active walkers joined the challenge.} &
\pcell{14}{``The leaderboard has 24 times as many users, so report 11{,}200.'' Respond.}{1.5cm}{No: size does not fix a non-chance method. Report 6{,}900.} \\ \hline
\end{probgrid}
\\ \hline
\end{guidednotes}

% ── THE NOTES END AT GUIDED PRACTICE ─────────────────────────────────────────
% There is deliberately no "you do" here: ../homework/ is the individual practice,
% started in class in the last 8 minutes while the teacher circulates.

\end{document}
